% Part: computability
% Chapter: recursive-functions
% Section: examples

\documentclass[../../../include/open-logic-section]{subfiles}

\begin{document}

\olfileid[pt]{cmp}{rec}{exa}
\olsection{Exemplos de funções recursivas primitivas}


Já temos alguns exemplos de funções recursivas primitivas: as funções adição
e multiplicação~$\Add$ e $\Mult$. A função identidade $\fn{id}(x) = x$ é
recursiva primitiva, pois é simplesmente~$\Proj{1}{0}$. As funções constantes
$\fn{const}_n(x) = n$ são recursivas primitivas, pois podem ser definidas a
partir de $\Zero$ e $\Succ$ por composições sucessivas. Isso é útil quando
queremos usar constantes em definições recursivas primitivas; p.ex., se
queremos definir a função $f(x) = 2 \cdot x$, podemos obtê-la por composição a
partir de $\fn{const}_n(x)$ e da multiplicação, tomando
$f(x) = \Mult(\fn{const}_2(x), \Proj{1}{0}(x))$. Daqui em diante, usaremos
esse recurso.

\begin{prop}
  A função exponenciação $\fn{exp}(x, y) = x^y$ é recursiva primitiva.
\end{prop}

\begin{proof}
  Podemos definir $\fn{exp}$ por recursão primitiva como
  \begin{align*}
    \fn{exp}(x, 0) & = 1\\
    \fn{exp}(x, y+1) & = \Mult(x, \fn{exp}(x,y)).
    \intertext{Estritamente falando, essa não é uma definição recursiva a
      partir de funções recursivas primitivas. Formalmente, porém, temos:}
    \fn{exp}(x, 0) & = f(x)\\
    \fn{exp}(x, y+1) & = g(x, y, \fn{exp}(x,y)).
    \intertext{onde}
    f(x) & = \Succ(\Zero(x)) = 1\\
    g(x, y, z) & = \Mult(\Proj{3}{0}(x, y, z), \Proj{3}{2}(x, y, z)) = x \cdot z
  \end{align*}
  e, portanto, $f$ e $g$ são definidas por composição a partir de funções
  recursivas primitivas.
\end{proof}

\begin{prop}
  A função predecessor $\fn{pred}(y)$ definida por
  \[
  \fn{pred}(y) = \begin{cases}
    0 & \text{se $y=0$}\\
    y-1 & \text{caso contrário}
  \end{cases}
  \]
  é recursiva primitiva.
\end{prop}

\begin{proof}
 Observe que
 \begin{align*}
   \fn{pred}(0) & = 0 \text{ e}\\
   \fn{pred}(y+1) & = y.
 \end{align*}
 Isso é quase uma definição por recursão primitiva. Estritamente falando, não
 se ajusta ao padrão de definição por recursão primitiva, pois esse padrão
 exige pelo menos um argumento adicional~$x$. Também é peculiar porque não
 usa de fato $\fn{pred}(y)$ na definição de $\fn{pred}(y+1)$. Mas podemos
 definir primeiro $\fn{pred}'(x, y)$ por
 \begin{align*}
   \fn{pred}'(x, 0) & = \Zero(x) = 0,\\
   \fn{pred}'(x, y+1) & = \Proj{3}{1}(x, y, \fn{pred'}(x, y)) = y.
 \end{align*}
e então definir $\fn{pred}$ a partir dela por composição, p.ex., como
$\fn{pred}(x) = \fn{pred}'(\Zero(x), \Proj{1}{0}(x))$.
\end{proof}

\begin{prop}
  A função fatorial $\fn{fac}(x) = \fact{x} = 1 \cdot 2 \cdot 3
  \cdot \dots \cdot x$ é recursiva primitiva.
\end{prop}

\begin{proof}
  A definição recursiva primitiva evidente é
  \begin{align*}
    \fn{fac}(0) &= 1\\
    \fn{fac}(y+1) & = \fn{fac}(y) \cdot (y+1).
    \intertext{Formalmente, temos primeiro de definir uma função $h$ de dois argumentos}
    h(x, 0) & = \fn{const}_1(x)\\
    h(x, y+1) & = g(x, y, h(x, y))
    \intertext{onde $g(x, y, z) = \Mult(\Proj{3}{2}(x, y, z),
      \Succ(\Proj{3}{1}(x, y, z)))$; depois, tomamos}
    \fn{fac}(y) & = h(\Proj{1}{0}(y), \Proj{1}{0}(y)) = h(y,y).
  \end{align*}
  Daqui em diante, seremos um pouco menos formais e não forneceremos as
  definições completas por composição e recursão primitiva.
\end{proof}

\begin{prop}
  A subtração truncada, $x \tsub y$, definida por
  \[
  x \tsub y = \begin{cases}
    0 & \text{se $x < y$}\\
    x-y & \text{caso contrário}
  \end{cases}
  \]
  é recursiva primitiva.
\end{prop}

\begin{proof}
  Temos:
  \begin{align*}
    x \tsub 0 & = x\\
    x \tsub (y+1) & = \fn{pred}(x \tsub y) 
  \end{align*}
\end{proof}

\begin{prop}
 A distância entre $x$ e $y$, $\left|x-y\right|$, é recursiva primitiva.
\end{prop}

\begin{proof}
  Temos $\left| x-y \right| = (x \tsub y) + (y \tsub x)$; portanto, a
  distância pode ser definida por composição a partir de $+$ e $\tsub$, que
  são recursivas primitivas.
\end{proof}

\begin{prop}
  O máximo de $x$ e $y$, $\fn{max}(x,y)$, é recursivo primitivo.
\end{prop}

\begin{proof}
  Podemos definir $\fn{max}(x,y)$ por composição a partir de $+$ e $\tsub$:
  \[
  \fn{max}(x,y) \defis x + (y \tsub x).
  \]
  Se $x$ é o máximo, O.s.\@, $x \ge y$, então $y \tsub x = 0$, de modo que
  $x + (y \tsub x) = x + 0 = x$. Se $y$ é o máximo, então
  $y \tsub x = y - x$ e, portanto,
  $x + (y \tsub x) = x + (y - x) = y$.
\end{proof}

\begin{prop}
  \ollabel{prop:min-pr}
  O mínimo de $x$ e $y$, $\fn{min}(x,y)$, é recursivo primitivo.
\end{prop}

\begin{proof}
  Exercício.
\end{proof}

\begin{prob}
  Demonstre \olref[cmp][rec][exa]{prop:min-pr}.
\end{prob}

\begin{prob}
Mostre que \[
f(x, y) =
2^{(2^{\iddots^{2^{x}}})}\raisebox{1ex}{\bigg\rbrace}
\raisebox{1ex}{\text {$y$ ocorrências de $2$}}\] é recursiva primitiva.
\end{prob}

\begin{prob}
Mostre que a divisão inteira $d(x, y) = \lfloor x/y \rfloor$ (O.s.\@, a
divisão na qual se despreza tudo o que vem depois da vírgula decimal) é
recursiva primitiva. Quando $y = 0$, estipulamos $d(x, y) = 0$. Dê uma
definição explícita de~$d$ usando recursão primitiva e composição.
\end{prob}

\begin{prop}
O conjunto das funções recursivas primitivas é fechado sob as duas operações
seguintes:
\begin{enumerate}
\item Somas finitas: se $f(\vec x, z)$ é recursiva primitiva, então também é
a função
\[
g(\vec x, y) \defis \sum_{z = 0}^y f(\vec x, z).
\]
\item Produtos finitos: se $f(\vec x, z)$ é recursiva primitiva, então também
é a função
\[
h(\vec x, y) \defis \prod_{z = 0}^y f(\vec x, z).
\]
\end{enumerate}
\end{prop}

\begin{proof}
Por exemplo, as somas finitas são definidas recursivamente pelas equações
\begin{align*}
  g(\vec x, 0) & = f(\vec x, 0)\\
  g(\vec x, y+1) & = g(\vec x, y) + f(\vec x, y+1).
\end{align*}
\end{proof}


\end{document}
